\section{OrdMon (Monoides Ordenados)}

Os limites na categoria \textbf{OrdMon} devem respeitar simultaneamente a estrutura algébrica do monoide e a relação de ordem parcial compatível.

\subsection{Fundamentação Teórica}
\begin{itemize}
    \item \textbf{Produtos:} O objeto $M \times N$ herda a ordem do produto, onde $(m, n) \le (m', n')$ se e só se as relações se verificam individualmente em cada componente ($m \le_M m'$ e $n \le_N n'$).
    \item \textbf{Equalizadores:} Correspondem ao sub-monoide formado pelos elementos onde os morfismos coincidem, mantendo a ordem herdada do domínio original.
\end{itemize}

\subsection{Implementação Computacional}
Abaixo detalham-se as funções para a criação do produto e a extração do equalizador.

\begin{lstlisting}[caption={Construções em Monoides Ordenados}]
% Produto em OrdMon
function OM = ordmon_product(M1, M2)
    OM.elements = set_product(M1.elements, M2.elements);
    OM.mul = @(a, b) [M1.mul(a(1), b(1)), M2.mul(a(2), b(2))];
    OM.unit = [M1.unit, M2.unit];
    % Verificacao da ordem combinada
    OM.le = @(a, b) M1.le(a(1), b(1)) && M2.le(a(2), b(2));
end

% Equalizador em OrdMon
function E = ordmon_equalizer(f, g, M)
    E.elements = {};
    for i = 1:length(M.elements)
        if isequal(f(M.elements{i}), g(M.elements{i}))
            E.elements{end+1} = M.elements{i};
        end
    end
    E.mul = M.mul;
    E.le = M.le;
end
\end{lstlisting}